\clearpage
\fancyhead[L]{\small INHALT.}
\SourceLeaf{667}
\begingroup
\raggedright
\noindent\hangindent=1em Connoissance des tems ou des mouvemens célestes, à l'usage des astronomes et des navigateurs, pour l'an 1810; publiée par le bureau des longitudes. De l'imprimerie impériale 1808. Paris \mbox{.\ .\ .\ .\ .\ .\ .\ .} \hfill\mbox{1809 Juni 17}\quad\mbox{Seite 515}\par\smallskip
\noindent\hangindent=1em \textsc{von Lindenau}: Resultate von Untersuchungen über den Durchmesser der Sonne \mbox{.\ .\ .\ .\ .\ .\ .\ .\ .\ .\ .\ .\ .\ .\ .\ .\ .\ .\ .\ .\ .\ .\ .\ .} \hfill\mbox{1809 Juli 17}\quad\mbox{--- 517}\par\smallskip
\noindent\hangindent=1em \textsc{De Rossel}: Voyage de \textsc{Dentrecasteaux}, envoyé à la recherche de la Pérouse. Publié par ordre de Sa Majesté l'Empereur et Roi. Tome second. Paris 1808 \mbox{.\ .\ .\ .\ .\ .\ .\ .\ .\ .\ .\ .\ .\ .\ .\ .\ .\ .\ .\ .\ .\ .\ .} \hfill\mbox{1809 Oct. 21}\quad\mbox{--- 520}\par\smallskip
\noindent\hangindent=1em \textsc{Regner}: Supplementum ad historiam de Parallaxeos Solaris inventione \mbox{.\ .} \hfill\mbox{1810 Jan. 29}\quad\mbox{--- 523}\par\smallskip
\noindent\hangindent=1em \textsc{von Lindenau}: [Greenwicher Sonnen-Beobachtungen von 1787---1798] \mbox{.\ .\ .} \hfill\mbox{1810 Juni 16}\quad\mbox{--- 525}\par\smallskip
\noindent\hangindent=1em Connoissance des tems ou des mouvemens célestes à l'usage des astronomes et des navigateurs, pour l'an 1811, publiée par le bureau des longitudes. Paris 1809 \mbox{.\ .\ .\ .\ .\ .\ .\ .\ .\ .\ .\ .\ .\ .\ .\ .\ .\ .\ .\ .\ .} \hfill\mbox{1810 Sept. 24}\quad\mbox{--- 327}\par\smallskip
\noindent\hangindent=1em \textsc{Knitlmayer}: Versuch einer genauen Darstellung des Progressionsverhältnisses der Planeten- und Trabantenabstände von ihren Centralkörpern. Brünn 1808 \mbox{.\ .\ .\ .\ .\ .\ .\ .\ .\ .\ .\ .\ .\ .\ .\ .\ .\ .\ .\ .\ .\ .} \hfill\mbox{1810 Nov. 10}\quad\mbox{--- 528}\par\smallskip
\noindent\hangindent=1em \textsc{Bessel}: Untersuchungen über die scheinbare und wahre Bahn des im Jahre 1807 erschienenen grossen Cometen. Königsberg \mbox{.\ .\ .\ .\ .\ .\ .\ .\ .} \hfill\mbox{1811 Febr. 16}\quad\mbox{--- 530}\par\smallskip
\noindent\hangindent=1em \textsc{Schulze}: Prediger in Polenz: Darstellung des Weltsystems. Ein Leitfaden für den Unterricht in der Astronomie auf Schulen. Leipzig 1811 \mbox{.\ .} \hfill\mbox{1811 Febr. 23}\quad\mbox{--- 533}\par\smallskip
\noindent\hangindent=1em \textsc{Seyffer}: Super longitudine geographica speculae astronomicae regiae, quae Monachii est, ex triginta septem defectionibus Solis observatis et ad calculos revocatis nunc primum definita. München 1810 \mbox{.\ .\ .\ .\ .} \hfill\mbox{1811 Mai 18}\quad\mbox{--- 533}\par\smallskip
\noindent\hangindent=1em \textsc{Nouet}: Description de l'Egypte. Tome premier à Paris de l'imprimerie impériale 1809. Mémoire premier. Observations astronomiques, faites en Egypte pendant les années VI, VII et VIII \mbox{.\ .\ .\ .\ .\ .\ .\ .\ .} \hfill\mbox{1811 Mai 18}\quad\mbox{--- 534}\par\smallskip
\noindent\hangindent=1em \textsc{Pasquich}: Epitome elementorum astronomiae sphaerico-calculatoriae.\par\noindent\hangindent=1em\hspace{1em}Pars prima: Elementa theoretića astronomiae sphaerico-calculatoriae.\par\noindent\hangindent=1em\hspace{1em}Pars secunda: Elementa practica astronomiae sphaerico-calculatoriae. Wien 1811 \mbox{.\ .\ .\ .\ .\ .\ .\ .\ .\ .\ .\ .\ .\ .\ .\ .\ .\ .\ .\ .\ .\ .\ .} \hfill\mbox{1811 Sept. 9}\quad\mbox{--- 535}\par\smallskip
\noindent\hangindent=1em \textsc{Klüber}: Die Sternwarte zu Mannheim. Mannheim, Heidelberg 1811 \mbox{.\ .\ .} \hfill\mbox{1811 Decbr. 2}\quad\mbox{--- 536}\par\smallskip
\noindent\hangindent=1em Connoissance des tems, etc.\ .\ .\ . pour l'an 1812. Paris 1810 \mbox{.\ .\ .\ .\ .\ .} \hfill\mbox{1811 Decbr. r4}\quad\mbox{--- 537}\par\smallskip
\noindent\hangindent=1em D'Aquila: Decouverte de l'orbite de la terre du point central de l'orbite du soleil, leur situation et leur forme; de la section du zodiaque, par le plan de l'équateur, et du mouvement concordant des deux globes. Paris 1806 \mbox{.\ .\ .\ .\ .\ .\ .\ .\ .\ .\ .\ .\ .\ .\ .\ .\ .\ .\ .\ .\ .\ .} \hfill\mbox{1812 Juni 27}\quad\mbox{--- 541}\par\smallskip
\noindent\hangindent=1em \textsc{Harding}: Himmelsatlass. Lieferung III. und IV. \mbox{.\ .\ .\ .\ .\ .\ .\ .\ .\ .} \hfill\mbox{1812 Octob. 3}\quad\mbox{--- 541}\par\smallskip
\noindent\hangindent=1em \textsc{Gerling}: Methodi projectionis orthographicae usum ad calculos parallacticos facilitandos explicavit, simulque eclipsin solarem die VII. Sept. 1820 apparituram hoc modo tractatam mappaque geographica illustratam tamquam exemplum proposuit. Göttingen \mbox{.\ .\ .\ .\ .\ .\ .\ .\ .\ .\ .} \hfill\mbox{1812 Nov. 28}\quad\mbox{--- 542}\par\smallskip
\noindent\hangindent=1em \textsc{Carlini}: Esposizione di un nuovo methodo di costruire le tavole astronomiche applicato alle tavole del sole. Mailand 1810. \mbox{.\ .\ .\ .\ .\ .\ .} \hfill\mbox{1812 Decbr. 3}\quad\mbox{--- 543}\par\smallskip
\noindent\hangindent=1em \textsc{De Zach}: Nouvelles tables d'aberration et de nutation pour quatorze cent \SourceLeaf{668}quatre étoiles, avec une table générale d'aberration pour les planètes et les comètes, précédées d'une instruction qui renferme l'explication de l'usage de ces tables, suivies de plusieurs nouvelles tables destinées à faciliter les calculs astronomiques. Marseille \mbox{.\ .\ .\ .\ .\ .\ .\ .\ .} \hfill\mbox{1812 Decbr. 7}\quad\mbox{Seite 546}\par\smallskip
\noindent\hangindent=1em \textsc{Bode}: Astronomisches Jahrbuch für das Jahr 1814. Berlin 1811\par\noindent\hangindent=1em\hspace{1em}Astronomisches Jahrbuch für das Jahr 1815. Berlin 1812 \mbox{.\ .\ .\ .\ .} \hfill\mbox{1813 Febr. 6}\quad\mbox{--- 548}\par\smallskip
\noindent\hangindent=1em Tables astronomiques publiées par le bureau des longitudes de France.\par\noindent\hangindent=1em \textsc{Burckhardt}: Tables de la lune. Paris 1812. \mbox{.\ .\ .\ .\ .\ .\ .\ .\ .\ .\ .} \hfill\mbox{1813 März 8}\quad\mbox{--- 553}\par\smallskip
\noindent\hangindent=1em Connaissance des tems etc.\ .\ .\ . pour l'an 1813. Paris 1811.\par\noindent\hangindent=1em Connaissance des tems etc.\ .\ .\ . pour l'an 1814. Paris 1812 \mbox{.\ .\ .\ .\ .\ .} \hfill\mbox{1813 Juni 3}\quad\mbox{--- 554}\par\smallskip
\noindent\hangindent=1em \textsc{Carlini} e \textsc{Brioschi}: Effemeridi astronomische di Milano per l'anno bisestile 1812. Milano 1811.\par\noindent\hangindent=1em\hspace{1em}Effemeridi astronomiche di Milano per l'anno 1813 \mbox{.\ .\ .\ .\ .\ .\ .\ .} \hfill\mbox{1813 Juni 14}\quad\mbox{--- 556}\par\smallskip
\noindent\hangindent=1em \textsc{De Lindenau}: Investigatio nova orbitae a Mercurio circa Solem descriptae: accedunt tabulae Planetae ex elementis recens correctis et ex theoria gravitatis ill. de la Place constructae \mbox{.\ .\ .\ .\ .\ .\ .\ .\ .\ .\ .\ .} \hfill\mbox{1813 Juli 5}\quad\mbox{--- 559}\par\smallskip
\noindent\hangindent=1em \textsc{Bessel}: Bestimmung der Praecession [Berliner Preisschrift] \mbox{.\ .\ .\ .\ .\ .} \hfill\mbox{1814 Jan. 31}\quad\mbox{--- 561}\par\smallskip
\noindent\hangindent=1em \textsc{Pond}: Beobachtung der Sommersolstitien von 1812 und 1813, und des Wintersolstitium 1812 auf der königlichen Sternwarte zu Greenwich. Philosophical Transactions of the Royal Society of London for the Jear 1813 \mbox{.\ .\ .\ .\ .\ .\ .\ .\ .\ .\ .\ .\ .\ .\ .\ .\ .\ .\ .\ .\ .\ .\ .} \hfill\mbox{1814 Febr. 14}\quad\mbox{--- 561}\par\smallskip
\noindent\hangindent=1em \textsc{Bode}: Astronomisches Jahrbuch für das Jahr 1816 \mbox{.\ .\ .\ .\ .\ .\ .\ .\ .\ .} \hfill\mbox{1814 März 10}\quad\mbox{--- 563}\par\smallskip
\noindent\hangindent=1em \textsc{Carlini}: Effemeridi astronomiche di Milano per l'anno 1814 \mbox{.\ .\ .\ .\ .\ .} \hfill\mbox{1814 März 24}\quad\mbox{--- 566}\par\smallskip
\noindent\hangindent=1em \textsc{Mollweide}: Commentationes mathematico-philologicae tres, sistentes explicationem duorum locorum difficilium, alterius Virgilii, alterius Platonis itemque examinationem duorum mensurarum praeceptorum Columellae. Leipzig 1813' \mbox{.\ .\ .\ .\ .\ .\ .\ .\ .\ .\ .\ .\ .\ .\ .\ .\ .\ .} \hfill\mbox{1814 Mai 2}\quad\mbox{--- 567}\par\smallskip
\noindent\hangindent=1em \textsc{De Zach}: L'attraction des montagnes, et ses effets sur les fils à plomb ou sur les niveaux des instrumens d'astronomie,\ .\ .\ .\ . Avignon 1814 \mbox{.\ .} \hfill\mbox{1814 Dec. 31}\quad\mbox{--- 569}\par\smallskip
\noindent\hangindent=1em \textsc{Carlini}: Effemeridi astronomiche di Milano per l'anno 1815 \mbox{.\ .\ .\ .\ .\ .} \hfill\mbox{1815 Jan. 19}\quad\mbox{--- 574}\par\smallskip
\noindent\hangindent=1em \textsc{Bode}: Astronomisches Jahrbuch für das Jahr 1817 \mbox{.\ .\ .\ .\ .\ .\ .\ .\ .\ .} \hfill\mbox{1815 Febr. 18}\quad\mbox{--- 576}\par\smallskip
\noindent\hangindent=1em Connaissance des tems etc.\ .\ .\ . pour l'an 1816 \mbox{.\ .\ .\ .\ .\ .\ .\ .\ .\ .\ .} \hfill\mbox{1815 März 11}\quad\mbox{--- 581}\par\smallskip
\noindent\hangindent=1em \textsc{Soldner}: Neue Methode, beobachtete Azimuthe zu reduciren \mbox{.\ .\ .\ .\ .} \hfill\mbox{1815 März 23}\quad\mbox{--- 587}\par\smallskip
\noindent\hangindent=1em \textsc{Stefenelli}: Beiträge zur Berechnung beobachteter Azimuthe \mbox{.\ .\ .\ .\ .\ .} \hfill\mbox{1815 März 23}\quad\mbox{--- 587}\par\smallskip
\noindent\hangindent=1em \textsc{Piazzi}: Praecipuarum stellarum inerrantium positiones etc. Palermo 1814 \mbox{.} \hfill\mbox{1815 Juni 10}\quad\mbox{--- 590}\par\smallskip
\noindent\hangindent=1em \textsc{Müller}: Berechnung der Sonnenfinsternisse \mbox{.\ .\ .\ .\ .\ .\ .\ .\ .\ .\ .\ .} \hfill\mbox{1816 Febr. 3}\quad\mbox{--- 591}\par\smallskip
\noindent\hangindent=1em \textsc{Carlini}: Effemeridi astronomiche di Milano per l'anno bisestile 1816 \mbox{.\ .\ .} \hfill\mbox{1816 März 23}\quad\mbox{--- 592}\par\smallskip
\noindent\hangindent=1em \textsc{Harding}: Himmelsatlas. Lieferung V. \mbox{.\ .\ .\ .\raisebox{.4ex}{\rotatebox{-60}{-}}\ .\ .\ .\ .\ .\ .\ .\ .\ .\ .} \hfill\mbox{1817 Jan. 9}\quad\mbox{--- 593}\par\smallskip
\noindent\hangindent=1em Connaissance des tems etc.\ .\ .\ . pour l'an 1818 \mbox{.\ .\ .\ .\ .\ .\ .\ .\ .\ .\ .} \hfill\mbox{1817 April 5}\quad\mbox{--- 595}\par\smallskip
\noindent\hangindent=1em \textsc{Schumacher}: De latitudine speculae Manhemiensis. Copenhagen 1816 \mbox{.\ .} \hfill\mbox{1817 Mai 12}\quad\mbox{--- 597}\par\smallskip
\noindent\hangindent=1em Tavole delle parallassi, di altezza e di latitudine calcolate dagli astronomi dell' osservatorio\ .\ .\ . nel collegio Romano. Rom 1816 \mbox{.\ .\ .\ .\ .\ .\ .} \hfill\mbox{1817 Juli 7}\quad\mbox{--- 599}\par\smallskip
\noindent\hangindent=1em \textsc{Carlini}: Effemeridi astronomiche di Milano per l'anno 1817. Mailand 1816 \mbox{.} \hfill\mbox{1817 Sept. 8}\quad\mbox{---\textperiodcentered{} 601}\par\smallskip
\endgroup
